This section is very similar to \Cref{sec:shift_pm_one_vector}.

We will show that, the rank-one shift family of vectors $x \mathbf{vv^T}$ is also $\NP$-complete to decide when $x\in\mathbb{R}$ and $\mathbf{v}\in\{0,1\}^n$. Let $S$ be the set of symmetric matrices and $\DD$ as in \Cref{def:diagonally_dominant}.

\begin{definition}[$\ROSP^{0,1}$]
\label{def:rank_one_zero_vector}
	\begin{equation}
		\ROSP^{0,1} := \left\{M \in \mathbb{R}^{n\times n} \cap S, 	n\in\mathbb{N} \;\vert\; \exists \mathbf{v}\in \{0,1\}^n, \exists x\in\mathbb{R} \text{ s.t. } (M - x \mathbf{vv}^T) \in \DD \right\}
	\end{equation}
\end{definition}

Our hope would be to be able to say that $\ROSP^{0,1}\in\P$, so that we can automatically improve on DeVille's $x1^{n\times n}$ shift \cite{deville_optimizing_2019}. Unfortunately, this is not the case:

\begin{theorem}
\label{thm:rosp_zero_one_np_complete}
	$\ROSP^{0,1}\in \NP$-complete.
\end{theorem}
\begin{proof}
	The proof, as much as in the case of \Cref{thm:rosp_pm_one_npc}, goes by reducing $\EThreeSAT\leq^p_m \ROSP^{0,1}$. The fact that $\ROSP^{0,1}$ is in $\NP$ is trivial \hubi{always assuming that $x$ admits a reasonable representation}.
	
	Given a formula $\phi$, we will denote with $n$ the number of variables $x_1, \dots, x_n$, and $m$ the number of clauses. We define $f =\max(m, 2n+10) - 1$. If $m < f + 1$, then we simply repeat the $m$-th clause $f - m + 1$ times, and call $m$ the new number of clauses. Of course, repeating a clause does not change the truth value of formula $\phi$.
	
	We will assume $\phi = C_1 \land C_2 \land \dots \land C_m$, i.e., a conjunction of 3-clauses, and we will use the shorthand notation $j\in C_i$ or $j\notin C_i$ to denote that variable $x_j$ is present/not in clause $C_i$ with any sign. We will use $x_j\in C_i$ or $\neg x_j\in C_i$ to denote the sign of variable $j$ in clause $C_i$.
	
	The reduction $R$ mapping $\EThreeSAT$ formula $\phi$ to matrix $R(\phi) = M^\phi\in\mathbb{R}^{n+f+1\times n + f + 1}$ is as follows:
	\begin{itemize}
		\item $M^\phi_{i,i} = 2 (n + f + 1)$, $\forall i \in \{1, \dots, n\}$.
		\item  $M^\phi_{i,j} = 0$, $\forall i\neq j \in \{1, \dots, n\}^2$.
		\item $M^\phi_{i,i} = \frac{(n-3)}{2} + 3$, $\forall i \in \{n+1, \dots, n+f+1\}. $
		\item   $M^\phi_{i,j} = 1$, $\forall i\neq j \in \{n+1, \dots, n+f+1\}^2$.
		\item $M^\phi_{i,j} = \begin{cases}
			\frac{1}{2} & j \notin C_{i-n} \\
			0 & \neg x_j \in C_{i-n} \\
			1 & x_j \in C_{i-n}
		\end{cases}$, $\forall (i, j) \in\{n+1, \dots, n+f+1\}\times \{1, \dots, n\}$.
		\item $M_{i,j}^\phi = M_{j,i}$, $\forall (i,j) \in\{1, \dots, n\}\times \{n+1, \dots, n+f+1\}$.
	\end{itemize}
	
	In the same spirit as in the proof of \Cref{thm:rosp_pm_one_npc}, we will show the following claims:
	
	\begin{claim}[$R(\phi)$ is symmetric]
	\label{claim:reduction_symmetric_rosp_zero_one}
		$\forall \phi \in \Phi, R(\phi)\in S$.
	\end{claim}
	\begin{proof}
		The claim is true by construction.
	\end{proof}
	
	\begin{claim}[Polynomial reduction]
	\label{claim:polynomial_reduction_rosp_zero_one}
		Computing $R(\phi)$ takes polynomial time in $n, m$. Hence, $R$ is a polynomial reduction.
	\end{claim}
	\begin{proof}
		The claim is also true by construction: it is just enough to fill in the matrix $M^\phi$ using the rules described above, every entry has constant/polylogarithmic representation size, and the size of the matrix is polynomial in $n, m$.
	\end{proof}
	
	\begin{claim}
	\label{claim:right_implication_rosp_zero_one}
		$\phi\in\EThreeSAT \implies R(\phi) \in \ROSP^{0,1}$.
	\end{claim}
	
	\begin{claim}
	\label{claim:left_implication_rosp_zero_one}
		$\phi\notin\EThreeSAT \implies R(\phi) \notin \ROSP^{0,1}$.
	\end{claim}
	
	The above \Cref{claim:right_implication_rosp_zero_one,claim:left_implication_rosp_zero_one} are going to be proved below.
	
	It is clear that, with \Cref{claim:reduction_symmetric_rosp_zero_one,claim:polynomial_reduction_rosp_zero_one,claim:right_implication_rosp_zero_one,claim:left_implication_rosp_zero_one} in place, \Cref{thm:rosp_zero_one_np_complete} is true.
\end{proof}

\begin{proof}[Proof of \Cref{claim:right_implication_rosp_zero_one}]
	Given a formula $\phi\in\EThreeSAT$ with true/false certificate $\mathbf{c}=[v_1, \dots, v_n]\in\{0,1\}^n$, it is enough to pick $x=1$ and $\mathbf{v} = \mathbf{c} \cup 1^{f+1}$ (where, with the $\cup$ operator, we mean a concatenation).
	
	To show that $(R(\phi) - 1 \mathbf{vv}^T) \in \DD$, it is enough to notice that:
	
	\begin{itemize}
		\item For rows $i\in\{1, \dots, n\}$, the diagonally dominant condition is: 
		\begin{align}
			M_{i,i}^\phi - v_i^2 - \sum_{j=1,j\neq i}^{n+f+1} |M_{i,j}^\phi - v_i v_j| & \geq \\
			2(n+f+1) - 1 - \sum_{j=1,j\neq i}^{n} |M_{i,j}^\phi - v_i v_j| - \sum_{j=n+1}^{n+f+1} |M_{i,j}^\phi - v_i v_j| & \geq \\
			2n + 2f + 1 - (n-1) - (f+1) & > 0
		\end{align}
		\item For rows $i\in\{n+1, \dots, f+n+1\}$, instead, the diagonally dominant condition is:
		\begin{align}
			M_{i,i}^\phi - v_i^2 - \sum_{j=1, j\neq i}^{n+f+1} |M_{i,j}^\phi - v_i v_j| & = \\
			\frac{n-3}{2}+3 - 1 - \sum_{j\leq n }|M_{i,j}^\phi - v_i v_j| - \sum_{j>n, j\neq i} |M_{i,j}^\phi - v_i v_j| & = \\
			\frac{n}{2} + \frac{1}{2} - \sum_{j\leq n, j\in C_{i-n}}|M_{i,j}^\phi - v_i v_j| - \sum_{j\leq n, j\notin C_{i-n}}|M_{i,j}^\phi - v_i v_j| & = \\
			\frac{n}{2} + \frac{1}{2} - \sum_{j\leq n, j\in C_{i-n}}|M_{i,j}^\phi - v_i v_j| - \frac{n-3}{2} & \geq \\
			2 - 2 & \geq 0 \label{eq:at_least_one_correct_sign}
		\end{align}
	\end{itemize}

Where \Cref{eq:at_least_one_correct_sign} is because, within clause $C_{i-n}$, the given valid certificate $\mathbf{c} = \{v_1, \dots, v_n\}$ has at least one (out of three) variable $v_j$ with the correct sign 0 or 1, whereas $v_i=1$ for $i > n$. So, the summation $\sum_{j\leq n, j\in C_{i-n}}|M_{i,j}^\phi - v_i v_j| \leq 2 $.

Since we proved that for all rows $i\in\{1, \dots, n+f+1\}$ the diagonally dominant condition of $R(\phi)$ is met, this concludes the proof.
\end{proof}

\begin{proof}[Proof of \Cref{claim:right_implication_rosp_zero_one}]
	Assume $\phi\notin\EThreeSAT$. First, we prove that when $x\neq1$, then no vector $\mathbf{v}\in\{0,1\}^{n+f+1}$ can make $(R(\phi) - x\mathbf{vv}^T)\in \DD$. Then, when we fix $x=1$, we show that for every shift $1 \mathbf{vv}^T$ there must exist a row $i>n$ s.t. the diagonally dominant condition is not met. This will conclude the proof that $\phi\notin\EThreeSAT \implies R(\phi)\notin\ROSP^{0,1}$.
	
	To begin with, it has to hold that $x\leq 1$: because all off-diagonal entries of $R(\phi) = M^\phi$ have norm $\leq 1$, selecting $x>1$ would worsen the diagonally dominant condition on every row. 
	
	\begin{fact} 
	\label{fact:shift_finish_with_ones}
		It must hold that $v_i=1$ for all $i\in\{n+1, \dots, n+f+1\}$; otherwise, the matrix $R(\phi) - x\mathbf{vv}^T$ would not meet the diagonally dominant condition on all rows $i>n$ s.t. $v_i=0$. Hence, from now on, we will assume $\mathbf{v}\in \{0,1\}^n\cup 1^{f+1}$.
	\end{fact}
	
	When $x\leq \frac{1}{2}$, the diagonally dominant condition on row $i>n$ is:
	\begin{align}
	M_{i,i}^\phi - x v_i^2 -  \sum_{j=1}^n|M_{i,j}^\phi - x v_j v_i| - \sum_{j=n+1, j\neq i}^{n+f+1}|M_{i,j}^\phi - x v_j v_i| & < \\
	\frac{n-3}{2} + 3 - x - 0 - f(1 - x) & \leq \\
	\frac{n-3}{2} + 3 - x  - (2n + 10)(1-x) & \leq \\
	\frac{n-3}{2} + 3 - n - 5 & < 0
	\end{align}
	
	Hence, the only feasible $x\in (\frac{1}{2}, 1]$.
	
	We now show that, also when $\frac{1}{2} < x < 1$, no $\mathbf{v}$ is a valid shift: if we interpret the first $n$ entries of any vector $\mathbf{v}\in\{0,1\}^n \cup 1^{f+1}$ as a certificate for $\phi\in\EThreeSAT$, it must hold that there is at least a clause $C_i = (s_{j_1} x_{j_1} \lor s_{j_2}x_{j_2} \lor s_{j_3}x_{j_3})$ ($x_{j_l}$ being variables, and $s_{j_l}$ the signs) where $v_{j_l} \neq s_{j_l}$ for all $l\in\{1, 2, 3\}$ (or, the formula $\phi$ would be satisfiable).
	
	For such a clause $i\in\{1,\dots,m\}$, the diagonally dominant condition on row $i+n$ is:
		
	\begin{align}
		M^\phi_{i+n,i+n} - x v_{i+n}^2 - \sum_{j = 1}^{n} |M^\phi_{i+n,j} - x v_{i+n} v_j| -  \sum_{j = n+1, j\neq i+n}^{n+f+1} |M^\phi_{i+n,j} - x v_{i+n} v_j| & = \\
		\frac{n-3}{2} + 3 - x -  \sum_{j = 1}^{n} |M^\phi_{i+n,j} - x v_{i+n} v_j| - f(1-x) & = \\
		\frac{n+3}{2} - x - f + fx -  \sum_{j=1,j\in C_i}^n |M^\phi_{i+n,j} - x v_{i+n} v_j| - \sum_{j=1, j\notin C_i}^n |M^\phi_{i+n,j} - x v_{i+n} v_j| & \leq \\
		\frac{n+3}{2} - x - f + fx - 3x - (x-\frac{1}{2})(n-3) & = \\
		\frac{n + 3 - 2x - 2f + 2fx - 6x - 2xn + n + 6x - 3}{2} & = \\
		\frac{2n - 2x - 2f + 2 fx - 2xn}{2} & < \\
		\frac{(1-x)(2n) - (1-x)(2f)}{2} & = \\
		\frac{(1-x)(2n-2f)}{2} & < 0
	\end{align}
	
We have just proved that unless $x=1$, no $\mathbf{v}\in\{0,1\}^{n}\cup 1^{f+1}$ (and, due to \Cref{fact:shift_finish_with_ones}, no $\mathbf{v}\in\{0,1\}^{n+f+1}$) can make $(R(\phi) - x\mathbf{vv}^T)\in \DD$.

We conclude the proof by showing that, even when $x=1$, for any vector $\mathbf{v}\in \{0,1\}^{n+f+1}$ the matrix  $(R(\phi) - x\mathbf{vv}^T)\notin \DD$.

Due to the above \Cref{fact:shift_finish_with_ones}, $\mathbf{v}\in\{0,1\}^n\cup 1^{f+1}$. Once again, if we consider the first $n$ entries of $\mathbf{v}$ as the truth assignment of variables $x_1, \dots, x_n$, we know there must be at least a clause $C_i$, $i\leq m$, s.t. all three variables have the wrong sign in the assignment $\mathbf{v}$. 

On row $i+n$, then, the diagonally dominant condition is:
\begin{align}
	M_{i+n, i+n}^\phi - v_{i+n}^2 - \sum_{j = 1}^n |M_{i+n, j}^\phi - v_j v_{i+n}| -  \sum_{j = n+1, j\neq i+n}^{f+n+1} |M_{i+n, j}^\phi - v_j v_{i+n}| & = \\
	\frac{n-3}{2} + 3 - 1 -  \sum_{j\leq n}|M_{i+n, j}^\phi - v_j v_{i+n}| & = \\
	\frac{n+1}{2} - \sum_{j=1, j \in C_i}^n |M_{i+n, j}^\phi - v_j v_{i+n}| - \sum_{j=1, j \notin C_i}^n |M_{i+n, j}^\phi - v_j v_{i+n}| & = \\
	\frac{n+1}{2} - 3 - \frac{n-3}{2} & = -1
\end{align}

So the diagonally dominant condition is not met.
This finally concludes the proof.

\end{proof}